% Expanded submission evidence. These appendices are intentionally excluded
% from the declared 8,000-word main-body count.

\newcommand{\EvidenceSheet}[7]{%
\clearpage
\section{#1}
\noindent
\begin{tabularx}{\textwidth}{>{\bfseries}p{3.5cm}X}
\toprule
Identifier & #2 \\
Purpose & #3 \\
Preconditions & #4 \\
Procedure or decision & #5 \\
Acceptance evidence & #6 \\
Residual risk & #7 \\
\bottomrule
\end{tabularx}
\vspace{0.8cm}

\noindent\textbf{Traceability note.}
This record connects a stated project requirement to a concrete implementation
or verification activity. It is retained as submission evidence so that the
system can be assessed against observable behaviour rather than feature claims
alone.
}

\chapter{Requirements Catalogue and Traceability}
\label{app:traceability}

This appendix expands the project objectives into assessable requirements.
Priority follows the MoSCoW convention: Must, Should, Could, and Won't for the
current release. The verification column identifies the strongest available
form of evidence.

\begin{longtable}{p{1.4cm}p{8.7cm}p{2.0cm}p{2.2cm}}
\caption{Functional requirements catalogue.}\label{tbl:functional_requirements}\\
\toprule
\textbf{ID} & \textbf{Requirement} & \textbf{Priority} & \textbf{Evidence} \\
\midrule
\endfirsthead
\toprule
\textbf{ID} & \textbf{Requirement} & \textbf{Priority} & \textbf{Evidence} \\
\midrule
\endhead
FR-01 & The application shall launch and expose its core destination catalogue without an active internet connection. & Must & Airplane-mode scenario \\
FR-02 & The application shall load all 300 curated Gandaki Province destination records from bundled assets. & Must & Dataset validation \\
FR-03 & A user shall be able to search and filter destinations by meaningful tourism attributes. & Must & Widget and manual tests \\
FR-04 & A destination detail view shall present descriptive, geographic, seasonal, activity, and image information. & Must & UI inspection \\
FR-05 & The map shall render bundled Gandaki tiles when online tile services are unavailable. & Must & Airplane-mode scenario \\
FR-06 & The routing service shall prefer a cached route, then an online OSRM route, then an offline Haversine fallback. & Must & Routing branch tests \\
FR-07 & Users shall be able to save destinations and retain them between application sessions. & Must & Persistence test \\
FR-08 & Users shall be able to submit and retain local destination reviews. & Should & Persistence test \\
FR-09 & The system shall produce explainable personalized destination recommendations. & Must & Scenario evaluation \\
FR-10 & Recommendation capability shall remain available without the backend. & Must & Offline recommendation test \\
FR-11 & The chatbot shall answer tourism-oriented questions using local intelligence when disconnected. & Must & Intelligence end-to-end test \\
FR-12 & Emergency-like messages shall receive deterministic safety-oriented handling. & Must & Safety scenario test \\
FR-13 & Translation shall support the application's declared English and Nepali travel phrases. & Should & Translation evaluation \\
FR-14 & The application shall queue user interactions locally when the backend is unavailable. & Should & Synchronisation test \\
FR-15 & Queued interactions shall be submitted to the backend after connectivity returns. & Should & Integration scenario \\
FR-16 & The backend shall expose authenticated registration, login, and current user endpoints. & Should & Backend tests \\
FR-17 & The backend shall expose destination and accommodation catalogue endpoints. & Must & API contract test \\
FR-18 & The backend shall expose recommendation, similar-destination, and evaluation endpoints. & Must & API contract test \\
FR-19 & Administrators shall be able to rebuild the semantic destination index. & Could & Endpoint inspection \\
FR-20 & The application shall communicate backend unavailability clearly without preventing local use. & Must & Offline UX inspection \\
\bottomrule
\end{longtable}

\clearpage
\begin{longtable}{p{1.4cm}p{8.7cm}p{2.0cm}p{2.2cm}}
\caption{Non-functional requirements catalogue.}\label{tbl:nonfunctional_requirements}\\
\toprule
\textbf{ID} & \textbf{Requirement} & \textbf{Priority} & \textbf{Evidence} \\
\midrule
\endfirsthead
\toprule
\textbf{ID} & \textbf{Requirement} & \textbf{Priority} & \textbf{Evidence} \\
\midrule
\endhead
NFR-01 & Core discovery, saved data, local recommendations, maps, and local assistance shall tolerate a complete loss of network connectivity. & Must & Airplane-mode suite \\
NFR-02 & The Android release shall remain practical to distribute despite bundled maps, images, and recommendation assets. & Must & APK size inspection \\
NFR-03 & Personal data and authentication secrets shall not be embedded in public application assets. & Must & Configuration review \\
NFR-04 & Backend passwords shall be stored using a one-way password hashing mechanism. & Must & Auth service test \\
NFR-05 & Recommendation results shall include understandable reasons rather than scores alone. & Should & Response inspection \\
NFR-06 & Ranking evaluation shall use repeatable relevance scenarios and standard information-retrieval metrics. & Must & Evaluation script \\
NFR-07 & The app shall use accessible contrast, readable labels, and touch targets suitable for repeated mobile use. & Should & UI inspection \\
NFR-08 & Failures in online enhancement services shall degrade to local capability rather than crash the application. & Must & Failure injection \\
NFR-09 & Local route cache entries shall expire after a defined freshness period. & Should & Cache inspection \\
NFR-10 & Source modules shall retain clear separation between presentation, domain logic, storage, and remote services. & Should & Architecture review \\
NFR-11 & Automated static analysis shall complete without reported Flutter issues. & Must & \texttt{flutter analyze} \\
NFR-12 & Backend automated tests shall pass in the documented environment. & Must & \texttt{pytest} result \\
NFR-13 & The application shall be maintainable through central configuration of its backend base URL. & Should & Configuration review \\
NFR-14 & Dataset fields shall use a consistent schema that can be validated before packaging. & Must & Dataset validation \\
NFR-15 & Safety responses shall be deterministic and shall not depend exclusively on an external language model. & Must & Emergency tests \\
\bottomrule
\end{longtable}

\EvidenceSheet
{Offline Launch and Core Discovery}
{REQ-TRACE-01 / FR-01, FR-02, NFR-01}
{Demonstrate that the application remains useful at initial launch when no network is available.}
{A release build is installed; airplane mode is enabled; bundled assets have not been removed.}
{Launch the app, dismiss onboarding if required, open the home destination feed, perform a destination search, and open a destination detail view. Repeat after force-closing the application.}
{The shell loads without a blocking network error; locally packaged destination content is visible; search and detail navigation remain available; the app can be relaunched.}
{Live events, weather, remote reviews, and newly changed destination information cannot be guaranteed while disconnected.}

\EvidenceSheet
{Bundled Destination Catalogue Integrity}
{REQ-TRACE-02 / FR-02, NFR-14}
{Establish that the packaged catalogue is structurally valid and sufficiently rich for filtering and recommendation.}
{The destination JSON asset is present and the application's destination model is available.}
{Parse the complete asset; check identifier uniqueness; validate coordinates; confirm required text, category, activity, season, budget, accessibility, family, culture, nature, and image fields; report malformed records.}
{All 300 intended destination records load through the real model mapping and can be used by discovery and recommendation services.}
{Curated data can become stale; a future release requires a governed update and validation process.}

\EvidenceSheet
{Offline Map Availability}
{REQ-TRACE-03 / FR-05, NFR-01}
{Confirm that map exploration does not disappear when an online tile provider cannot be reached.}
{The Gandaki MBTiles asset is packaged; location permission is granted or a known map location is selected.}
{Enable airplane mode, open the map, pan across Pokhara and surrounding Gandaki areas, change supported zoom levels, and reopen the map after restarting the app.}
{Raster tiles render from the bundled database for the declared coverage and zoom range without a tile-network dependency.}
{Coverage is deliberately bounded; street-level tiles outside the packaged region or zoom range are unavailable.}

\EvidenceSheet
{Three-Tier Route Resolution}
{REQ-TRACE-04 / FR-06, NFR-08, NFR-09}
{Verify deterministic routing behaviour across cached, online, and fully disconnected states.}
{Two valid coordinates are selected and the route cache can be inspected or cleared.}
{Request a route once online, repeat to exercise the cache, then clear the relevant entry and request the route in airplane mode. Compare route source, distance, and presentation.}
{The first eligible online request uses OSRM, the repeated request can use the SQLite cache, and the disconnected cache miss produces a Haversine fallback instead of a crash.}
{The Haversine line is an estimate and must not be presented as road-aware turn-by-turn navigation.}

\EvidenceSheet
{Persistent Saved Destinations and Reviews}
{REQ-TRACE-05 / FR-07, FR-08}
{Confirm that user-created local state survives navigation and process restarts.}
{At least one destination is available; application storage has not been manually cleared.}
{Save a destination, submit a local review, navigate away, force-close the app, relaunch, and inspect saved and reviewed content. Then remove the saved item and confirm the change persists.}
{SQLite-backed saved and review state remains consistent across sessions and reflects create and delete actions.}
{Device uninstall or storage clearing removes local-only state unless a future synchronisation feature has uploaded it.}

\EvidenceSheet
{Offline Recommendation Continuity}
{REQ-TRACE-06 / FR-09, FR-10, NFR-05}
{Demonstrate that personalization is not exclusively dependent on the FastAPI service.}
{Destination assets and local recommendation resources are packaged; user preferences can be supplied.}
{Generate recommendations online, disable connectivity, change a preference or context, and generate recommendations again. Inspect ordering, reasons, and absence of blocking errors.}
{The offline pipeline returns ranked destinations with human-readable reasons and uses locally available semantic and contextual signals.}
{Offline ranking does not have access to fresh cross-user interaction data or newly trained backend models.}

\EvidenceSheet
{Local Tourism Assistance and Safety}
{REQ-TRACE-07 / FR-11, FR-12, NFR-15}
{Verify that local intelligence can answer supported tourism questions and prioritise emergency handling.}
{Offline knowledge assets and intelligence services are initialized; network connectivity is disabled.}
{Ask for a destination suggestion, a practical travel phrase, a follow-up question requiring context, and an emergency-style message. Inspect response mode and safety content.}
{Supported tourism queries receive local retrieval or template responses; conversation state is retained; emergency content is handled by deterministic safety logic.}
{The assistant is not a substitute for emergency services, medical advice, or verified real-time trail conditions.}

\EvidenceSheet
{Backend Failure Degradation}
{REQ-TRACE-08 / FR-20, NFR-08, NFR-13}
{Confirm that a backend outage is visible but does not disable local application features.}
{The app is configured with a deliberately unreachable backend URL while local assets remain intact.}
{Launch the app, trigger a backend-enhanced action, inspect the user-facing state, then continue to discovery, map, saved items, local recommendations, and local chatbot functions.}
{The application communicates that online services are unavailable and continues to expose its offline-capable workflows.}
{Remote-only capabilities such as authentication, fresh analytics, and online model responses remain unavailable until connectivity and backend health recover.}

\chapter{Data Dictionary and Dataset Evidence}
\label{app:data_dictionary}

The recommendation and discovery experience depends on consistent structured
data. The following dictionary documents the principal destination fields used
across Flutter and FastAPI components.

\begin{longtable}{p{3.2cm}p{2.2cm}p{6.1cm}p{2.5cm}}
\caption{Destination data dictionary.}\label{tbl:data_dictionary}\\
\toprule
\textbf{Field} & \textbf{Type} & \textbf{Meaning and use} & \textbf{Validation} \\
\midrule
\endfirsthead
\toprule
\textbf{Field} & \textbf{Type} & \textbf{Meaning and use} & \textbf{Validation} \\
\midrule
\endhead
\texttt{id} & string & Stable destination key used by UI, persistence, routes, interactions, and recommendation outputs. & Required; unique \\
\texttt{name} & string & Human-readable destination name displayed throughout the application. & Required; non-empty \\
\texttt{description} & string & Curated descriptive content for details, retrieval, and semantic embedding. & Required; non-empty \\
\texttt{latitude} & number & Geographic latitude used by maps, distance features, and routing. & -90 to 90 \\
\texttt{longitude} & number & Geographic longitude used by maps, distance features, and routing. & -180 to 180 \\
\texttt{category} & string & Primary tourism category such as village, lake, viewpoint, temple, or trek. & Controlled value \\
\texttt{activities} & list & Supported experiences such as trekking, culture, nature, photography, or pilgrimage. & Non-empty list \\
\texttt{best\_season} & list & Seasons in which the destination is most suitable. & Controlled values \\
\texttt{budget\_level} & string & Relative cost category used by filters and contextual scoring. & Controlled value \\
\texttt{accessibility} & string & Relative access difficulty used in suitability filtering and explanations. & Controlled value \\
\texttt{adventure\_\allowbreak level} & integer & Ordinal indicator of desired physical/adventure intensity. & Bounded scale \\
\texttt{family\_\allowbreak friendly} & boolean & Indicates suitability for family-oriented discovery scenarios. & Boolean \\
\texttt{cultural\_\allowbreak richness} & number & Relative cultural value used in filters and ranking features. & Bounded scale \\
\texttt{nature\_score} & number & Relative natural-attraction value used in filters and ranking features. & Bounded scale \\
\texttt{tags} & list & Additional searchable and semantically useful descriptive labels. & Normalised strings \\
\texttt{image} & string & Local or resolvable image reference used by cards and detail presentation. & Asset existence \\
\texttt{sbert\_text} & string & Prepared semantic text combining meaningful destination attributes. & Non-empty text \\
\texttt{province} & string & Administrative region used to constrain and describe coverage. & Gandaki scope \\
\bottomrule
\end{longtable}

\EvidenceSheet
{Dataset Identity and Uniqueness}
{DATA-EVIDENCE-01}
{Prevent ambiguous persistence, interaction, and recommendation references.}
{The complete destination asset is available to a validation script or model loader.}
{Collect every destination identifier and name; report empty values, duplicate identifiers, and suspicious duplicate normalized names; verify that referenced identifiers remain stable between app and backend assets.}
{Every packaged destination has one stable unique identifier and a usable display name.}
{Name changes and future record merges require migration rules so saved items and interactions remain associated correctly.}

\EvidenceSheet
{Geographic Validity}
{DATA-EVIDENCE-02}
{Ensure that destination points can be rendered, measured, and routed without invalid coordinates.}
{All records have been parsed into numeric latitude and longitude values.}
{Check global coordinate bounds, inspect whether points fall within the intended Nepal/Gandaki coverage, and plot outliers for manual review.}
{All intended destinations produce valid map points and no record silently defaults to an unrelated coordinate.}
{A valid coordinate does not guarantee a safe or publicly accessible arrival point; field verification remains necessary.}

\EvidenceSheet
{Semantic Retrieval Readiness}
{DATA-EVIDENCE-03}
{Confirm that recommendation and chatbot retrieval operate on informative text rather than destination names alone.}
{Each record exposes description, category, activity, tag, and prepared semantic text fields.}
{Inspect representative records across categories; compare source attributes with prepared semantic text; check for empty or duplicated semantic content; verify embedding-index alignment with identifiers.}
{Semantic records contain enough distinct tourism context to support similarity search and scenario-oriented ranking.}
{Curated descriptions may reflect author bias and can underrepresent less documented communities.}

\EvidenceSheet
{Local Media Integrity}
{DATA-EVIDENCE-04}
{Verify that offline destination presentation can resolve packaged images reliably.}
{The Flutter asset manifest and local image directory are available.}
{Resolve each destination image reference, identify missing or duplicate assets, open a sample across image sizes, and inspect card/detail fallback behaviour.}
{Intended local image references resolve and missing media does not crash a destination card or details screen.}
{Image licensing, accessibility text, and future content updates require an explicit governance process.}

\EvidenceSheet
{Accommodation Relationship Integrity}
{DATA-EVIDENCE-05}
{Confirm that accommodation records can be associated with their destination context.}
{The accommodation dataset and destination identifier set are available.}
{Validate accommodation identifiers, destination foreign-key references, coordinates, names, types, price data, and duplicated records; inspect destinations with zero and many linked accommodations.}
{Each of the 771 intended accommodation records is parseable and every declared destination relationship resolves.}
{Availability and prices are not real-time and must be labelled accordingly in future booking-oriented releases.}

\chapter{Detailed API Contract Catalogue}
\label{app:api_contracts}

Each contract sheet records the behaviour expected by the Flutter client and by
administrative or evaluation tooling. Exact payloads remain versioned by the
backend implementation; the sheets focus on externally observable obligations.

\EvidenceSheet
{Health Check Contract}
{API-01 / \texttt{GET /health}}
{Provide a low-cost signal that the FastAPI process is reachable before online enhancement is attempted.}
{The backend process has started and its configured port is reachable.}
{Send an unauthenticated GET request; inspect status code and response body; repeat while a downstream optional model provider is unavailable.}
{The endpoint returns a successful service-health response without requiring user credentials or an external model provider.}
{Process health alone does not prove that every optional dependency or data source is healthy.}

\EvidenceSheet
{Destination Collection Contract}
{API-02 / \texttt{GET /destinations}}
{Expose the destination catalogue for online consumers and consistency checks.}
{The backend destination repository has loaded its dataset.}
{Request the collection; inspect record count, identifiers, names, coordinates, and pagination or limit behaviour if configured.}
{The response is valid JSON, uses stable identifiers, and exposes fields compatible with the client destination model.}
{Large future datasets require explicit pagination and version-aware client handling.}

\EvidenceSheet
{Destination Detail and Accommodation Contract}
{API-03 / \texttt{GET /destinations/\{id\}} and accommodation route}
{Expose one destination and its linked accommodation context.}
{A valid destination identifier and an invalid test identifier are available.}
{Request a valid detail record, request its accommodations, then request a missing identifier and inspect error semantics.}
{Valid requests return the intended record and relationships; missing identifiers return a controlled not-found response rather than a server error.}
{Accommodation data remains informational unless integrated with a verified availability provider.}

\EvidenceSheet
{Recommendation Contract}
{API-04 / \texttt{POST /recommend}}
{Return ranked and explainable destination suggestions for a supplied user and context.}
{The destination index and recommendation service are initialized; a valid request payload is available.}
{Submit a cold-start request, a preference-rich request, and a request with contextual attributes; inspect ordering, scores, reasons, metadata, and repeatability.}
{The endpoint returns a bounded ranked list with stable identifiers, interpretable reasons, and metadata sufficient for client presentation and analysis.}
{Sparse or synthetic interactions limit confidence in collaborative and learned ranking signals.}

\EvidenceSheet
{Similar Destination Contract}
{API-05 / \texttt{GET /similar/\{id\}}}
{Support related-destination discovery using semantic similarity.}
{A valid indexed destination identifier is available.}
{Request similar destinations, verify that the source destination is excluded, inspect similarity ordering, and test a missing identifier.}
{The response contains valid alternatives with scores or ordering consistent with the semantic index.}
{Semantic similarity can reinforce catalogue-description bias and does not automatically provide diversity.}

\EvidenceSheet
{Interaction Logging Contract}
{API-06 / \texttt{POST /interactions} and \texttt{/interactions/batch}}
{Accept behavioural events used for analytics and future recommendation improvement.}
{The interaction repository is writable and test user/destination identifiers are available.}
{Submit a single event, submit a mixed batch, inspect accepted counts, test validation failures, and verify repository persistence.}
{Valid events are stored once with required context; invalid events receive controlled validation responses; batch results expose partial-failure semantics.}
{Interaction collection requires consent, retention controls, and protection against event duplication or manipulation.}

\EvidenceSheet
{Chat Contract}
{API-07 / \texttt{POST /chat}}
{Provide online-enhanced tourism assistance while preserving controlled fallback behaviour.}
{The chat route is registered; at least one provider or fallback service is available.}
{Send representative discovery, route, translation, follow-up, unsupported, and provider-failure prompts; inspect response source and safety handling.}
{The endpoint returns a structured response and does not expose provider exceptions directly to the client.}
{External provider content can be inaccurate; the client must preserve safety disclaimers and local fallback.}

\EvidenceSheet
{Offline Chat Contract}
{API-08 / \texttt{POST /chat/offline}}
{Expose the backend's deterministic offline-oriented retrieval response for explicit testing and fallback use.}
{Offline chat data and service initialization have completed.}
{Submit supported and unsupported tourism prompts and inspect intent, response text, and destination references.}
{The endpoint responds without calling a cloud language model and returns controlled content for supported local knowledge.}
{The available response range is narrower than generative online assistance.}

\EvidenceSheet
{Translation Contract}
{API-09 / \texttt{POST /translate}}
{Provide supported travel translation with explicit language-pair and failure semantics.}
{Translation resources are initialized and a supported-pair matrix is documented.}
{Submit representative supported phrases, mixed-case input, romanised Nepali, an unsupported pair, and empty or excessive input.}
{Supported requests return structured translated text and language metadata; invalid requests receive controlled validation errors.}
{Coverage remains domain-specific and should not be represented as general professional translation.}

\EvidenceSheet
{Authentication Contract}
{API-10 / registration, login, and current-user routes}
{Provide identity for optional synchronized and personalized online capability.}
{The authentication repository and secret configuration are available.}
{Register a new account, reject duplicate or invalid input, login with correct and incorrect credentials, then access the current-user route with valid, expired, and missing tokens.}
{Passwords are not returned or stored in plain text; successful login returns a usable token; protected access rejects invalid credentials.}
{Authentication security also depends on deployment TLS, secret rotation, rate limiting, and production monitoring.}

\chapter{Verification Test Catalogue}
\label{app:test_catalogue}

The following cases describe repeatable acceptance and regression checks. They
supplement automated unit and widget tests with integrated offline-first
scenarios that cross storage, presentation, networking, and intelligence
boundaries.

\EvidenceSheet
{Cold Offline Startup Test}
{TEST-ACC-01}
{Test the most important failure condition: no connection at the start of a session.}
{Install the release build; clear recent tasks; enable airplane mode before launching.}
{Launch, complete or skip onboarding, open home, search, details, map, saved items, recommendations, chatbot, and translation in sequence.}
{Core local screens remain reachable; no infinite loading state or uncaught network exception prevents continued use.}
{Some screens can correctly show that optional online enhancement is unavailable.}

\EvidenceSheet
{Warm Online-to-Offline Transition Test}
{TEST-ACC-02}
{Check graceful degradation when connectivity disappears during active use.}
{The application is open and the backend is initially reachable.}
{Load online recommendations and a road route, disable connectivity, repeat both actions, then open local content and send a chatbot query.}
{Cached or local alternatives are used where designed; the state change is communicated without a crash.}
{An in-flight remote request may take until its configured timeout before fallback occurs.}

\EvidenceSheet
{Offline-to-Online Recovery Test}
{TEST-ACC-03}
{Check that the app recovers online capability without requiring reinstallation or data loss.}
{The app has been used offline and contains queued interactions.}
{Restore connectivity, trigger health checking or retry, request an online route, and observe interaction synchronization.}
{Online capability returns; queued events can be submitted; local saved state and caches remain intact.}
{Recovery timing depends on platform connectivity signals and backend availability.}

\EvidenceSheet
{Route Cache Precedence Test}
{TEST-ACC-04}
{Prove that an eligible cached route takes precedence and avoids unnecessary network use.}
{A known OSRM route has been cached and remains within the configured freshness period.}
{Request the identical coordinate pair with network available and inspect route source and outgoing traffic; repeat after expiry or cache removal.}
{The fresh cache serves the repeat request; expired or absent entries allow later routing tiers to execute.}
{Coordinate rounding and cache-key design can affect hit rate for nearly identical points.}

\EvidenceSheet
{Haversine Fallback Test}
{TEST-ACC-05}
{Confirm that a disconnected cache miss still returns a useful directional estimate.}
{No relevant route is cached and the device is offline.}
{Request a route between two valid destinations and compare reported straight-line distance with an independently calculated Haversine value.}
{A clearly identified fallback route is returned with plausible distance and without turn-by-turn claims.}
{Terrain, roads, bridges, closures, and walking feasibility are not represented by the fallback line.}

\EvidenceSheet
{Recommendation Scenario Evaluation Test}
{TEST-REC-01}
{Measure ranked retrieval quality against the report's five defined tourism scenarios.}
{Scenario queries, relevance sets, and the candidate destination catalogue are fixed for the run.}
{Generate top-ten rankings for every scenario and calculate Precision@10, Recall@10, nDCG@10, and MRR using the same relevance judgments.}
{Metric calculation is reproducible and results can be compared with the documented hybrid recommender values.}
{Five scenarios are useful for project evaluation but insufficient to establish population-level effectiveness.}

\EvidenceSheet
{Recommendation Explanation Test}
{TEST-REC-02}
{Verify that recommendation reasons correspond to actual ranking features and destination attributes.}
{A set of recommendations with score components and destination records is available.}
{Inspect reasons for season, activity, budget, distance, culture, nature, popularity, and similar-item signals; compare each statement with source attributes.}
{Reasons are understandable and do not claim attributes that the destination record or ranking path does not support.}
{Simplified explanations cannot fully describe every interaction within a learned ranking model.}

\EvidenceSheet
{Cold-Start Recommendation Test}
{TEST-REC-03}
{Confirm that a new user receives useful, diverse suggestions without an interaction history.}
{Use a new anonymous identifier or clear local recommendation history.}
{Request recommendations with no preferences, then with onboarding preferences; inspect relevance, diversity, and explanation changes.}
{The system supplies reasonable defaults and responds to explicit preferences without requiring collaborative history.}
{Popularity-based defaults can overexpose already prominent destinations.}

\EvidenceSheet
{Emergency Detection Test}
{TEST-AI-01}
{Verify deterministic priority handling of messages that indicate immediate danger or need for assistance.}
{The local intelligence pipeline is initialized and external AI providers are disabled.}
{Submit representative emergency, injury, lost-person, and benign messages in supported language forms; inspect detection and response precedence.}
{Emergency-like messages trigger the safety layer before normal tourism recommendation generation; benign queries avoid false emergency escalation.}
{Language variation and ambiguous messages require continuing review with domain experts.}

\EvidenceSheet
{Dialogue Continuity Test}
{TEST-AI-02}
{Check that follow-up questions can use recent conversation context.}
{Start a fresh conversation and ask for a destination recommendation with explicit constraints.}
{Send follow-ups that omit the destination name but refer to cost, route, season, and alternatives; then reset the conversation.}
{Relevant recent slots and entities are retained before reset and are not incorrectly reused after reset.}
{Long conversations and topic changes can still create ambiguity requiring clarification.}

\EvidenceSheet
{Translation Coverage Test}
{TEST-AI-03}
{Evaluate supported travel phrases, normalization, and controlled fallback behaviour.}
{The local glossary, phrasebook, templates, and available online adapter are known.}
{Test exact phrases, punctuation and case variants, romanised Nepali, Devanagari input, unsupported text, and provider unavailability.}
{Supported local phrases translate consistently; response metadata distinguishes translation source or fallback state.}
{General prose, specialist terminology, and dialect variation remain outside validated coverage.}

\EvidenceSheet
{Static Analysis and Automated Regression Test}
{TEST-ENG-01}
{Provide repeatable engineering-quality evidence for Flutter and backend components.}
{Dependencies are installed in the documented environments.}
{Run Flutter static analysis, focused intelligence and recommendation tests, backend pytest, and relevant widget tests; retain command, date, environment, and outcome.}
{Static analysis and focused suites produce a clear pass/fail signal; known failures are documented rather than concealed.}
{Environment-specific timing and golden-image rendering can differ across machines and SDK versions.}

\chapter{User Manual and Operational Workflows}
\label{app:user_manual}

This manual describes the intended end-user workflows. Screen names may evolve,
but the operational sequence and offline expectations remain part of the
project's acceptance baseline.

\EvidenceSheet
{First Launch and Onboarding}
{USER-01}
{Help a first-time traveler configure useful preferences without blocking immediate exploration.}
{The application is installed and launched for the first time.}
{Review onboarding, choose available preferences where offered, grant only required permissions, and continue to the main dashboard.}
{The user reaches the application shell; declared preferences can influence discovery and recommendation; denied optional permissions do not prevent catalogue browsing.}
{Location-dependent convenience is reduced when location permission is denied.}

\EvidenceSheet
{Browse and Search Destinations}
{USER-02}
{Find relevant destinations from the bundled Gandaki catalogue.}
{The home or discovery screen is open.}
{Scroll destination content, enter a search term, apply available filters, clear the query, and open a result.}
{Results respond to the query and filters; destination cards remain readable; clearing controls restores broader discovery.}
{Search quality depends on the vocabulary represented in the curated dataset.}

\EvidenceSheet
{Read Destination Details}
{USER-03}
{Inspect practical and descriptive information before deciding to travel.}
{A destination card or recommendation has been selected.}
{Review description, images, category, activities, season, budget, accessibility, location, and available accommodation context; follow map or save actions.}
{Information is presented consistently and available local assets remain visible without a connection.}
{Static descriptions cannot replace local safety advice or real-time access information.}

\EvidenceSheet
{Save and Review a Destination}
{USER-04}
{Retain personal travel interests and feedback locally.}
{A destination detail view is open and local storage is available.}
{Toggle the saved state, locate the destination in saved content, create or update a review, restart the app, and inspect persistence.}
{Saved and review state remains available after restart and can be intentionally changed or removed.}
{Local-only state is lost if application storage is cleared or the app is uninstalled.}

\EvidenceSheet
{Explore the Offline Map}
{USER-05}
{Use geographic context in places with unreliable mobile data.}
{The map screen is available and the bundled tile database is packaged.}
{Open the map, move to a known Gandaki area, pan and zoom within supported coverage, select a destination marker, and inspect the details action.}
{The map remains usable within packaged coverage and destination points connect to discovery workflows.}
{The bundled raster map is a snapshot and does not show live closures or every street-level detail.}

\EvidenceSheet
{Request Directions}
{USER-06}
{Obtain the best available route representation for the current connectivity state.}
{Origin and destination coordinates are available; location permission is granted if current position is used.}
{Request a route; inspect whether it is cached, road-aware online, or straight-line fallback; review distance and any warning; repeat offline.}
{The route source is understandable and the app does not imply road guidance when only a Haversine estimate is available.}
{Users must verify local conditions and should not rely solely on the fallback for safety-critical navigation.}

\EvidenceSheet
{Receive Personalized Recommendations}
{USER-07}
{Discover destinations that align with interests, season, budget, accessibility, and context.}
{The recommendation view is open and local destination data has loaded.}
{Provide or change available preference controls, request recommendations, read explanation text, open a suggested destination, and repeat without connectivity.}
{Ranked results update meaningfully and expose reasons that help the user judge suitability.}
{Recommendations assist discovery but do not guarantee satisfaction, safety, or availability.}

\EvidenceSheet
{Use the Tourism Chatbot}
{USER-08}
{Ask natural-language tourism questions and receive the best available local or online-enhanced response.}
{The chatbot screen is open.}
{Ask a specific question, use a follow-up, request an alternative, disconnect the network, and ask another supported local question.}
{The conversation remains usable; response state distinguishes fallback where appropriate; local questions continue to receive controlled assistance.}
{Generated or retrieved answers can be incomplete and should be verified for high-stakes decisions.}

\EvidenceSheet
{Translate a Travel Phrase}
{USER-09}
{Support practical communication using available English and Nepali resources.}
{The translation screen is open and source/target languages can be selected.}
{Enter a supported travel phrase, translate it, swap direction where supported, test romanised or Devanagari input, and repeat offline.}
{Supported content produces a readable result and offline resources remain available for their declared coverage.}
{The feature is not a certified interpreter and should not be relied on for legal or medical communication.}

\EvidenceSheet
{Recognise Offline and Recovery States}
{USER-10}
{Help users understand which capabilities remain local and which require the network.}
{The application is running and connectivity can be toggled.}
{Disable the network, observe messaging while using local features, attempt an online enhancement, restore connectivity, and retry.}
{The application communicates backend or connectivity loss without hiding local capability and recovers online actions when service returns.}
{Connectivity indicators can lag behind the actual reachability of a particular external service.}

\chapter{Deployment, Configuration, and Troubleshooting}
\label{app:operations}

\EvidenceSheet
{Flutter Development Setup}
{OPS-01}
{Provide a repeatable setup path for application development and verification.}
{Flutter, Dart, Android tooling, and a supported device or emulator are available.}
{Open the \texttt{app} directory, obtain dependencies, inspect environment configuration, run static analysis and focused tests, then launch a debug build.}
{Dependencies resolve, analysis completes, tests provide a clear result, and the app launches against the intended local or remote backend configuration.}
{SDK and plugin upgrades can change build behaviour; versions should be recorded for reproducibility.}

\EvidenceSheet
{Backend Development Setup}
{OPS-02}
{Provide a repeatable FastAPI setup for recommendation, chat, authentication, and analytics development.}
{A supported Python environment and project dependencies are available.}
{Configure environment variables, initialize required data, run backend tests, start the API server, and request the health endpoint.}
{Tests and health checking provide a clear operational signal and registered routes are visible through FastAPI documentation.}
{Optional provider keys and production databases require separate secure configuration.}

\EvidenceSheet
{Backend URL and Environment Switching}
{OPS-03}
{Avoid scattered host literals and support emulator, device, and deployed environments.}
{The central application backend configuration is available.}
{Set the intended backend environment or base URL, launch the app, inspect the health request, and verify that feature services derive their URIs from the central configuration.}
{Runtime requests use the selected central base URL and backend failure is surfaced through the intended offline UX.}
{Physical devices, emulators, and remote deployments have different network addressing and firewall constraints.}

\EvidenceSheet
{Android Release Build and Asset Verification}
{OPS-04}
{Produce a distributable package containing the offline assets required by the project.}
{Release signing and Flutter Android tooling are configured.}
{Build the release APK, inspect size, install it on a clean device, enable airplane mode, and verify destination data, images, MBTiles, and local intelligence assets.}
{The release installs and core offline workflows operate without relying on development-machine files.}
{Bundled assets increase download size and must be revalidated after asset-manifest changes.}

\EvidenceSheet
{Operational Troubleshooting Guide}
{OPS-05}
{Provide a systematic response to common failures without masking data loss or unsafe behaviour.}
{Logs, app settings, and health endpoints are available to the developer or tester.}
{For startup failure, inspect asset parsing and storage initialization; for backend failure, inspect base URL and health; for blank maps, inspect MBTiles packaging; for stale routes, inspect cache age; for AI failure, inspect local initialization before provider status.}
{Troubleshooting isolates the affected tier and preserves local functionality wherever possible.}
{Clearing application data is a last resort because it removes saved local state and can hide the original fault.}

\chapter{Architecture Decision Records}
\label{app:adrs}

\EvidenceSheet
{Use Flutter for the Mobile Client}
{ADR-01 / Accepted}
{Deliver a maintainable Android application with a consistent component model and future cross-platform potential.}
{The project requires a rich mobile interface, local storage, maps, and broad plugin support.}
{Adopt Flutter and Dart; structure features around presentation, domain models, and services; use platform integrations only where required.}
{The implemented client provides the intended screens, local data access, maps, and testable Dart services from one codebase.}
{Plugin quality and native build complexity can affect platform-specific behaviour.}

\EvidenceSheet
{Treat Offline Operation as the Default Architecture}
{ADR-02 / Accepted}
{Preserve useful travel support in the low-connectivity conditions central to the research problem.}
{Core catalogue, media, maps, routing fallback, storage, recommendations, and assistance can be packaged or computed locally.}
{Design local capability first and layer remote enhancement behind reachability checks, timeouts, caching, queues, and controlled fallback.}
{The application remains useful during backend or network loss and does not place every workflow behind authentication.}
{Offline assets are static and increase package size; synchronization requires careful conflict and freshness handling.}

\EvidenceSheet
{Use FastAPI for Online Enhancement Services}
{ADR-03 / Accepted}
{Expose typed Python services for recommendation, semantic search, evaluation, chat, analytics, and authentication.}
{The recommendation and NLP ecosystem is primarily Python-based and benefits from automatic API documentation.}
{Implement versioned route modules and application services in FastAPI; keep offline-critical capability in Flutter.}
{Online services can evolve independently while the client preserves local fallback.}
{A production deployment still requires TLS, observability, scaling, rate limiting, and secret management.}

\EvidenceSheet
{Combine SQLite, Hive, and SharedPreferences}
{ADR-04 / Accepted}
{Match local persistence mechanisms to relational records, high-frequency cached objects, queues, and lightweight settings.}
{Saved destinations, reviews, routes, recommendation caches, interactions, and preferences have different access patterns.}
{Use SQLite for relational saved/review/route data, Hive for recommendation cache and pending interactions, and SharedPreferences for small settings.}
{Each local data category has a storage mechanism appropriate to its structure and lifecycle.}
{Multiple stores increase migration, backup, and debugging complexity.}

\EvidenceSheet
{Bundle MBTiles for Offline Map Coverage}
{ADR-05 / Accepted}
{Ensure that geographic context remains available without runtime tile downloads.}
{A bounded Gandaki coverage area and practical zoom range are acceptable for the project.}
{Package a local raster MBTiles database and use online services only as optional enhancement.}
{The app renders the intended Gandaki map region in airplane mode.}
{Higher zoom levels and expanded coverage would substantially increase package size.}

\EvidenceSheet
{Use Cache, OSRM, then Haversine for Routing}
{ADR-06 / Accepted}
{Provide the best available route response while guaranteeing a disconnected fallback.}
{Road-aware offline routing data is too large or complex for the current scope.}
{Resolve routes through a fresh SQLite cache, then online OSRM, then a clearly identified straight-line Haversine approximation.}
{Repeated routes are efficient, online routes can be road-aware, and disconnected cache misses still return useful orientation.}
{The fallback must never be confused with safe road or trail instructions.}

\EvidenceSheet
{Hybrid Explainable Recommendation Pipeline}
{ADR-07 / Accepted}
{Balance semantic relevance, context, interaction evidence, diversity, and understandable user-facing reasons.}
{The catalogue contains rich text and attributes, while real-user interaction volume remains limited.}
{Combine semantic candidates, collaborative/popularity signals, contextual reranking, optional learned rescoring, diversity control, and explanation generation.}
{Scenario evaluation demonstrates strong first-result placement and exposes measurable recall limitations.}
{Learned and collaborative components require representative consented usage data before strong production claims.}

\EvidenceSheet
{Keep Deterministic Local Safety and Intelligence}
{ADR-08 / Accepted}
{Avoid making tourism assistance and emergency handling wholly dependent on a network or generative provider.}
{Local language processing, retrieval, dialogue state, phrase resources, and emergency rules can run on-device.}
{Run deterministic safety handling and local intelligence for every relevant message, using online language models only as an optional enhancement.}
{Supported queries and emergency-like messages retain controlled handling during provider or network failure.}
{Local coverage is finite and requires continued multilingual testing and expert review.}

\chapter{Selected Implementation Listings}
\label{app:listings}

The following listings are drawn directly from the implemented project. They
are included as technical evidence rather than as a claim that every source
file belongs in the dissertation. Line numbers support discussion and review.

\clearpage
\section{FastAPI Recommendation Route}
\lstinputlisting[
  language=Python,
  numbers=left,
  basicstyle=\ttfamily\scriptsize,
  caption={Implemented recommendation API route.}
]{../backend/api/v1/recommend.py}

\clearpage
\section{Backend Recommendation Application Service}
\lstinputlisting[
  language=Python,
  numbers=left,
  basicstyle=\ttfamily\scriptsize,
  caption={Implemented hybrid recommendation application service.}
]{../backend/application/services/recommendation_service.py}

\clearpage
\section{Backend Contextual Reranker Excerpt}
\lstinputlisting[
  language=Python,
  firstline=1,
  lastline=260,
  numbers=left,
  basicstyle=\ttfamily\scriptsize,
  caption={Selected contextual reranking implementation.}
]{../backend/infrastructure/ml/contextual_reranker.py}

\clearpage
\section{Flutter Routing Service}
\lstinputlisting[
  numbers=left,
  basicstyle=\ttfamily\scriptsize,
  caption={Implemented Flutter routing service with fallback tiers.}
]{../app/lib/core/navigation/routing_service.dart}

\clearpage
\section{Offline Semantic Encoder}
\lstinputlisting[
  numbers=left,
  basicstyle=\ttfamily\scriptsize,
  caption={Implemented offline semantic encoding support.}
]{../app/lib/features/recommendations/data/services/offline_semantic_encoder.dart}

\clearpage
\section{Intelligence Orchestrator Excerpt}
\lstinputlisting[
  firstline=1,
  lastline=320,
  numbers=left,
  basicstyle=\ttfamily\scriptsize,
  caption={Selected orchestration logic for local and online intelligence.}
]{../app/lib/features/intelligence/services/intelligence_orchestrator.dart}

\clearpage
\section{Backend Collaborative Filtering Test}
\lstinputlisting[
  language=Python,
  numbers=left,
  basicstyle=\ttfamily\scriptsize,
  caption={Automated collaborative filtering tests.}
]{../backend/tests/test_collaborative_filter.py}

\clearpage
\section{Flutter Intelligence End-to-End Test}
\lstinputlisting[
  numbers=left,
  basicstyle=\ttfamily\scriptsize,
  caption={Automated local intelligence end-to-end tests.}
]{../app/test/intelligence/end_to_end_test.dart}

\chapter{Project Governance, Risk, and Evidence Index}
\label{app:governance}

\begin{longtable}{p{1.2cm}p{4.7cm}p{1.8cm}p{1.8cm}p{3.9cm}}
\caption{Project risk register.}\label{tbl:risk_register}\\
\toprule
\textbf{ID} & \textbf{Risk} & \textbf{Prob.} & \textbf{Impact} & \textbf{Control} \\
\midrule
\endfirsthead
\toprule
\textbf{ID} & \textbf{Risk} & \textbf{Prob.} & \textbf{Impact} & \textbf{Control} \\
\midrule
\endhead
R-01 & Rural connectivity prevents online maps, routes, or AI responses. & High & High & Local assets, caches, route fallback, and local intelligence. \\
R-02 & Curated destination information becomes stale. & Medium & High & Versioned dataset, validation, update governance, and visible limitations. \\
R-03 & Haversine fallback is misunderstood as a safe road route. & Medium & High & Explicit source labelling and warning language. \\
R-04 & Recommendation results reinforce popularity bias. & Medium & Medium & Diversity control, explanations, coverage metrics, and future real-user evaluation. \\
R-05 & External language model provides inaccurate advice. & Medium & High & Local deterministic safety layer, controlled fallback, and disclaimers. \\
R-06 & Packaged maps and images make the application too large. & Medium & Medium & Bounded coverage, WebP compression, and release-size inspection. \\
R-07 & Interaction logging creates privacy concerns. & Medium & High & Consent, minimisation, retention policy, pseudonymous identifiers, and secure deployment. \\
R-08 & Authentication or secret configuration is deployed insecurely. & Low & High & Password hashing, environment secrets, TLS, rotation, and production review. \\
R-09 & Local databases become inconsistent across app upgrades. & Medium & Medium & Schema versioning, migrations, tests, and cautious storage clearing. \\
R-10 & Automated tests differ across SDK or operating-system versions. & Medium & Medium & Record environment, use focused suites, and document known failures. \\
R-11 & Destination images have unclear rights or attribution. & Medium & High & Maintain provenance and licensing records before public release. \\
R-12 & Rural communities are represented inaccurately or without benefit. & Medium & High & Community review, respectful descriptions, corrections, and local partnership. \\
\bottomrule
\end{longtable}

\clearpage
\section{Ethical and Responsible-Use Checklist}
\begin{longtable}{p{5.0cm}p{8.8cm}}
\toprule
\textbf{Consideration} & \textbf{Project response and future obligation} \\
\midrule
Connectivity equity & Core travel information is packaged locally so lack of mobile data does not remove the entire service. \\
Community representation & Descriptions should be reviewed with local stakeholders and corrected when they oversimplify culture or accessibility. \\
Recommendation fairness & Popularity and interaction signals should be monitored for exposure imbalance and should not erase lesser-known destinations. \\
Privacy & Interaction data should be minimized, consented, protected, retained only as needed, and never silently treated as harmless. \\
Safety & Route fallbacks and chatbot answers must expose limitations and avoid implying verified emergency or medical authority. \\
Accessibility & Interfaces should support readable contrast, adequate touch targets, meaningful labels, and continued testing with diverse users. \\
Image rights & Every distributable image requires provenance, permission or licence compatibility, and appropriate attribution where required. \\
Environmental impact & Recommendations should avoid encouraging unsafe overcrowding or environmentally damaging access without local guidance. \\
Transparency & Users should understand when content is local, cached, online-enhanced, generated, approximate, or potentially stale. \\
Correction process & A future public deployment requires a straightforward way for communities and users to report inaccurate or harmful content. \\
\bottomrule
\end{longtable}

\clearpage
\section{Project Milestone Record}
\begin{longtable}{p{2.9cm}p{3.5cm}p{4.8cm}p{2.4cm}}
\toprule
\textbf{Phase} & \textbf{Primary output} & \textbf{Completion evidence} & \textbf{Review focus} \\
\midrule
Problem definition & Research question, scope, objectives & Proposal and literature review & Relevance and feasibility \\
Data preparation & Destination and accommodation assets & Validated structured datasets and local media & Coverage and consistency \\
Architecture & Offline-first design and service boundaries & Diagrams, ADRs, and implementation structure & Failure tolerance \\
Core Flutter app & Discovery, details, persistence, map & Running screens and local storage & Usability and offline operation \\
Routing & Cache, OSRM, and Haversine tiers & Routing service and integrated scenarios & Source clarity and safety \\
Recommender & Semantic, context, collaborative, ranking, explanation & Services, scenario metrics, and tests & Relevance and limitations \\
Intelligence & Local NLP, retrieval, dialogue, safety, translation & Focused tests and integrated UI & Reliability and safety \\
Backend & API routes, auth, interactions, analytics & FastAPI implementation and pytest & Contracts and security \\
Evaluation & Metrics, regression evidence, critical review & Results chapter and appendices & Honest interpretation \\
Submission & Publication-quality report and reproducible artifacts & Compiled PDF, bibliography, and evidence index & Submission quality \\
\bottomrule
\end{longtable}

\clearpage
\section{Submission Evidence Index}
\begin{longtable}{p{4.5cm}p{5.0cm}p{4.5cm}}
\toprule
\textbf{Evidence category} & \textbf{Location in report} & \textbf{Assessment value} \\
\midrule
Research problem and motivation & Introduction and literature review & Establishes significance and academic context \\
Requirements and scope & Objectives, scope, Appendices E--F & Enables traceable assessment against intended behaviour \\
Architecture & Design chapter, diagrams, Appendix L & Explains structural decisions and trade-offs \\
Implementation & Implementation chapter and Appendix M listings & Demonstrates concrete engineering work \\
Destination and accommodation data & Data sections and Appendix G & Supports catalogue, retrieval, and recommendation claims \\
Recommendation evaluation & Evaluation chapter and Appendix I & Provides repeatable quantitative quality evidence \\
Offline verification & Evaluation chapter and test catalogue & Tests the central research contribution \\
API contracts & Appendix H & Documents client-backend boundaries and failure semantics \\
User operation & Appendix J & Demonstrates complete intended workflows \\
Deployment and reproducibility & Appendix K & Supports repeatable setup, build, and troubleshooting \\
Risk and ethics & Critical evaluation and Appendix N & Demonstrates responsible limitations and future obligations \\
\bottomrule
\end{longtable}
